% Evans second edition, section 1.2, printed pp. 3--6; PDF pages 20--23.
% Chinese: printed pp. 2--5; PDF pages 20--23.
\section{Examples}
\label{sec:1-2}

There is no general theory known concerning the solvability of all partial
differential equations. Such a theory is extremely unlikely to exist, given the
rich variety of physical, geometric, and probabilistic phenomena which can be
modeled by PDE. Instead, research focuses on various particular partial
differential equations that are important for applications within and outside
of mathematics, with the hope that insight from the origins of these PDE can
give clues as to their solutions.

Following is a list of many specific partial differential equations of interest
in current research. This listing is intended merely to familiarize the reader
with the names and forms of various famous PDE. To display most clearly the
mathematical structure of these equations, we have mostly set relevant physical
constants to unity. We will later discuss the origin and interpretation of many
of these PDE.

Throughout $x\in U$, where $U$ is an open subset of $\R^n$, and $t\geq0$.
Also $Du=D_xu=(u_{x_1},\ldots,u_{x_n})$ denotes the gradient of $u$ with
respect to the spatial variable $x=(x_1,\ldots,x_n)$. The variable $t$ always
denotes time.

\subsection{Single partial differential equations}
\subsubsection*{a. Linear equations}
\begin{enumerate}
\item Laplace's equation
\[ \Delta u=\sum_{i=1}^n u_{x_ix_i}=0. \]
\item Helmholtz's (or eigenvalue) equation
\[ -\Delta u=\lambda u. \]
\item Linear transport equation
\[ u_t+\sum_{i=1}^n b^iu_{x_i}=0. \]
\item Liouville's equation
\[ u_t+\sum_{i=1}^n(b^iu)_{x_i}=0. \]
\end{enumerate}
\par\Needspace{65mm}
\PDEAnnotation{随流不变与质量守恒}{输运方程沿 $\dot X=b(X,t)$ 保持 $u(X(t),t)$ 不变。
Liouville 方程展开后多出 $(\operatorname{div}b)u$，描述密度随体积伸缩而变。
只有 $\operatorname{div}b=0$ 时两式一致。}
\begin{enumerate}[resume]
\item Heat (or diffusion) equation
\[ u_t-\Delta u=0. \]
\item Schrödinger's equation
\[ \mathrm i u_t+\Delta u=0. \]
\item Kolmogorov's equation
\[ u_t-\sum_{i,j=1}^n a^{ij}u_{x_ix_j}+\sum_{i=1}^n b^iu_{x_i}=0. \]
\item Fokker--Planck equation
\[ u_t-\sum_{i,j=1}^n(a^{ij}u)_{x_ix_j}-\sum_{i=1}^n(b^iu)_{x_i}=0. \]
\end{enumerate}
\par\Needspace{65mm}
\PDEAnnotation{算子与形式伴随}{令 $L=a^{ij}\partial_{ij}-b^i\partial_i$（重复指标求和）。
分部积分给出 $L^*\rho=\partial_{ij}(a^{ij}\rho)+\partial_i(b^i\rho)$。
两式分别为 $u_t=Lu$ 与 $\rho_t=L^*\rho$，漂移项符号因此不同。}
\begin{enumerate}[resume]
\item Wave equation
\[ u_{tt}-\Delta u=0. \]
\item Klein--Gordon equation
\[ u_{tt}-\Delta u+m^2u=0. \]
\item Telegraph equation
\[ u_{tt}+2du_t-u_{xx}=0. \]
\par\Needspace{65mm}
\PDEAnnotation{阶数与时间初值}{热方程、波方程都是二阶 PDE，但前者时间一阶，
后者时间二阶，通常分别需给 $u(0)$ 与 $(u(0),u_t(0))$。
梁方程总阶数为四、时间阶数为二；总阶数不能直接决定初值个数。}
\item General wave equation
\[ u_{tt}-\sum_{i,j=1}^n a^{ij}u_{x_ix_j}+\sum_{i=1}^n b^iu_{x_i}=0. \]
\item Airy's equation
\[ u_t+u_{xxx}=0. \]
\item Beam equation
\[ u_{tt}+u_{xxxx}=0. \]
\end{enumerate}

\subsubsection*{b. Nonlinear equations}
\begin{enumerate}
\item Eikonal equation
\[ |Du|=1. \]
\item Nonlinear Poisson equation
\[ -\Delta u=f(u). \]
\item $p$-Laplacian equation
\[ \operatorname{div}(|Du|^{p-2}Du)=0. \]
\item Minimal surface equation
\[ \operatorname{div}\left(\frac{Du}{(1+|Du|^2)^{1/2}}\right)=0. \]
\item Monge--Ampère equation
\[ \det(D^2u)=f. \]
\end{enumerate}
\par\Needspace{65mm}
\PDEAnnotation{退化与奇异}{对 $p$-Laplacian，$p>2$ 时主部在 $Du=0$ 退化；
$1<p<2$ 时其线性化系数在那里奇异。散度形式及相应弱形式仍有意义，
不能把只在 $Du\ne0$ 合法的展开直接用于临界点。}
\begin{enumerate}[resume]
\item Hamilton--Jacobi equation
\[ u_t+H(Du,x)=0. \]
\item Scalar conservation law
\[ u_t+\operatorname{div}\mathbf F(u)=0. \]
\item Inviscid Burgers' equation
\[ u_t+uu_x=0. \]
\item Scalar reaction-diffusion equation
\[ u_t-\Delta u=f(u). \]
\item Porous medium equation
\[ u_t-\Delta(u^\gamma)=0. \]
\end{enumerate}
\par\Needspace{65mm}
\PDEAnnotation{多孔介质的扩散系数}{在 $u>0$ 的光滑区，
$\Delta(u^\gamma)=\gamma u^{\gamma-1}\Delta u+
\gamma(\gamma-1)u^{\gamma-2}|Du|^2$。
通常取 $\gamma>1,u\geq0$，主部在 $u=0$ 退化；$0<\gamma<1$ 对应快扩散情形。}
\begin{enumerate}[resume]
\item Nonlinear wave equation
\[ u_{tt}-\Delta u+f(u)=0. \]
\item Korteweg--deVries (KdV) equation
\[ u_t+uu_x+u_{xxx}=0. \]
\item Nonlinear Schrödinger equation
\[ \mathrm i u_t+\Delta u=f(|u|^2)u. \]
\end{enumerate}

\subsection{Systems of partial differential equations}
\subsubsection*{a. Linear systems}
\begin{enumerate}
\item Equilibrium equations of linear elasticity
\[ \mu\Delta\mathbf u+(\lambda+\mu)D(\operatorname{div}\mathbf u)=\mathbf0. \]
\item Evolution equations of linear elasticity
\[ \mathbf u_{tt}-\mu\Delta\mathbf u
 -(\lambda+\mu)D(\operatorname{div}\mathbf u)=\mathbf0. \]
\item Maxwell's equations
\[
 \begin{cases}
 \mathbf E_t=\operatorname{curl}\mathbf B,\\
 \mathbf B_t=-\operatorname{curl}\mathbf E,\\
 \operatorname{div}\mathbf B=\operatorname{div}\mathbf E=0.
 \end{cases}
\]
\end{enumerate}
\par\Needspace{65mm}
\PDEAnnotation{约束的传播}{真空 Maxwell 系统中，
$\partial_t\operatorname{div}\mathbf E=0$、
$\partial_t\operatorname{div}\mathbf B=0$，因为散度作用于旋度为零。
散度条件是初值必须满足且由演化保持的约束，不能视作任意额外数据。}

\subsubsection*{b. Nonlinear systems}
\begin{enumerate}
\item System of conservation laws
\[ \mathbf u_t+\operatorname{div}\mathbf F(\mathbf u)=\mathbf0. \]
\item Reaction-diffusion system
\[ \mathbf u_t-\Delta\mathbf u=\mathbf f(\mathbf u). \]
\item Euler's equations for incompressible, inviscid flow
\[
 \begin{cases}
 \mathbf u_t+\mathbf u\cdot D\mathbf u=-Dp,\\
 \operatorname{div}\mathbf u=0.
 \end{cases}
\]
\item Navier--Stokes equations for incompressible, viscous flow
\[
 \begin{cases}
 \mathbf u_t+\mathbf u\cdot D\mathbf u-\Delta\mathbf u=-Dp,\\
 \operatorname{div}\mathbf u=0.
 \end{cases}
\]
\end{enumerate}
\par\Needspace{65mm}
\PDEAnnotation{压力是约束乘子}{对光滑不可压流取散度，得到
$-\Delta p=\sum_{i,j}(\partial_i u^j)(\partial_j u^i)$。
压力由速度及边界条件耦合确定；没有独立的压力演化方程。
仅看速度的扩散项会漏掉这一非局部约束。}

See Zwillinger [Zw] for a much more extensive listing of interesting PDE.


\PDEBodyDiagram{1.2}{Three evolution mechanisms.}
{Dashed curves are initial profiles and solid curves are later profiles. Transport translates a shape; heat spreads a Gaussian while preserving its integral; the one-dimensional wave equation with zero initial velocity splits a profile into two half-amplitude translates. The panel axes are schematic and use separate scales.}
{ch01-models}

\subsection*{Editorial supplement: two principal-part calculations}
For $1<p<\infty$ and a $C^2$ function at a point where $Du\ne0$, the
chain rule gives
\begin{align*}
 \operatorname{div}(|Du|^{p-2}Du)
 &=\sum_{i,j=1}^n A_{ij}(Du)u_{x_ix_j},\\
 A(q)&=|q|^{p-2}\left(I+(p-2)\frac{q\otimes q}{|q|^2}\right).
\end{align*}
The eigenvalue in the direction of $q$ is $(p-1)|q|^{p-2}$; in every
orthogonal direction it is $|q|^{p-2}$. This proves ellipticity at $q\ne0$
and explains its failure to be uniform as $|q|\to0$ or $|q|\to\infty$
when $p\ne2$. At $p=2$ the equation is exactly Laplace's equation.

Similarly, the minimal surface equation expands as
\[
 \frac1{\sqrt{1+|Du|^2}}
 \sum_{i,j=1}^n\left(\delta_{ij}-\frac{u_{x_i}u_{x_j}}{1+|Du|^2}\right)
 u_{x_ix_j}=0.
\]
The matrix in parentheses has eigenvalues $1$ transverse to $Du$ and
$(1+|Du|^2)^{-1}$ parallel to $Du$. Thus it is positive definite for each
finite gradient, and a bound on $|Du|$ supplies a uniform ellipticity bound.
Both calculations exhibit quasilinear second-order equations.
\par\Needspace{65mm}
\PDEAnnotation{Monge--Ampère 的分支}{当 $n\geq2$，行列式非线性地依赖 Hessian。
在 $D^2u$ 正定的凸解分支上，线性化系数为正定的余子式矩阵。
离开该分支，椭圆性没有保证；仅写 $\det D^2u=f$ 还不足以指定椭圆问题。}
